%% toc: P15 164 164.12
\section{Systematic power in Hempel and Oppenheim: proofs}
\label{sec:annex-p15-w02c-p15-029-hempel-oppenheim}

The notation is that of \cref{sec:p15-w02c-p15-029-hempel-oppenheim}.
Three facts from \cref{sec:p15-systematic-power} are used throughout.
First, \(D_v\) is true at exactly the assignment \(v\) and \(M_v\) is
false at exactly \(v\). Second, for singular sentences, agreement at every
ground assignment is first-order equivalence, and \(U\entails V\) holds
exactly when every assignment satisfying \(U\) satisfies \(V\) (the
ground-freedom argument of \cref{annex:p15.power-normal-forms}). Third,
every range measure is invariant under logical equivalence, monotone
under entailment and satisfies \(r(\neg U)=1-r(U)\)
(\cref{thm:p15-power-range-laws}).

\begin{annexitem}{annex:p15-w02c-p15-029-ratio}{thm:p15-w02c-p15-029-ratio}{The bounded fraction and the derived-to-initial ratio}
\begin{proof}
Write \(k=|K'|\) and \(j=|K_T|\). Since \(S\) is not valid, \(K'\) is
nonempty, so \(k\ge1\); and \(0\le j\le k\). By
\cref{thm:p15-power-optimal-split}, \(s_0(T,K)=j/k\).

If \(j<k\), then \(s_0<1\) and
\[
 \frac{s_0}{1-s_0}=\frac{j/k}{(k-j)/k}=\frac{j}{k-j}=q(T,K).
\]
If \(j=k\), then \(s_0=1\) and \(q=+\infty\) by definition.

Put \(\phi(x)=x/(1-x)\) for \(0\le x<1\) and \(\phi(1)=+\infty\).
For \(0\le x<y<1\),
\[
 \phi(y)-\phi(x)=\frac{y-x}{(1-x)(1-y)}>0,
\]
and \(\phi(x)<+\infty=\phi(1)\) for every \(x<1\). So \(\phi\) is strictly
increasing on \([0,1]\) with values in the extended reals. Since
\(q(T,K)=\phi(s_0(T,K))\) for every \(T\), for any \(T,T'\) we have
\(s_0(T,K)\le s_0(T',K)\) if and only if
\(\phi(s_0(T,K))\le\phi(s_0(T',K))\).

For the one-name example, \cref{thm:p15-power-source-examples} gives
\(k=3\) and \(j=1\), so \(q=1/2\). For each \(K_j\) of that theorem,
\(k=252\) and the number of minimal sentences entailed by \(T\) is
\(175\), so \(q=175/(252-175)=175/77=25/11\).
\end{proof}
\end{annexitem}

\begin{annexitem}{annex:p15-w02c-p15-029-renaming}{thm:p15-w02c-p15-029-renaming}{Renaming the vocabulary}
\begin{proof}
\emph{Transport of structures.} Let \(\mathfrak M'\) be a structure for
\(L'\). Define a structure \(\mathfrak M=\mathfrak M'\sigma\) for \(L\)
with the same domain, interpreting each constant \(c\) as
\(\sigma(c)^{\mathfrak M'}\) and each predicate \(P\) as
\(\sigma(P)^{\mathfrak M'}\). These interpretations have the right arity
because \(\sigma\) preserves arity. An induction on formulas, with the
same variable assignment on both sides, shows
\(\mathfrak M'\sigma\models U\) if and only if \(\mathfrak M'\models\sigma U\).
Atomic formulas hold on both sides by the definition of
\(\mathfrak M'\sigma\). The connective steps are immediate. The
quantifier steps hold because the two structures have the same domain.
Conversely, every structure \(\mathfrak M\) for \(L\) equals
\(\mathfrak M'\sigma\) for the \(L'\)-structure \(\mathfrak M'\) that
interprets \(\sigma(c)\) as \(c^{\mathfrak M}\) and \(\sigma(P)\) as
\(P^{\mathfrak M}\). This structure is well defined because \(\sigma\)
is a bijection.

\emph{Item 1.} If \(U\entails V\) and \(\mathfrak M'\models\sigma U\), then
\(\mathfrak M'\sigma\models U\), hence \(\mathfrak M'\sigma\models V\),
hence \(\mathfrak M'\models\sigma V\). Conversely, if
\(\sigma U\entails\sigma V\) and \(\mathfrak M\models U\), write
\(\mathfrak M=\mathfrak M'\sigma\). Then \(\mathfrak M'\models\sigma U\),
so \(\mathfrak M'\models\sigma V\) and \(\mathfrak M\models V\). Validity
is entailment from \(\top\), and \(\sigma\top=\top\). An atom
\(P(c_1,\ldots,c_k)\) goes to the atom
\(\sigma(P)(\sigma(c_1),\ldots,\sigma(c_k))\). Since \(\sigma\) is
bijective on constants and on predicates of each arity, this map is a
bijection of ground atoms.

\emph{Item 2.} Number the ground atoms of \(L\) as \(A_1,\ldots,A_n\)
and those of \(L'\) as \(A'_i=\sigma A_i\). Then \(\sigma D_v=D'_v\) and
\(\sigma M_v=M'_v\) for each assignment \(v\). Let \(K\) have
conjunction \(S\), so that \(\sigma K\) has conjunction \(\sigma S\).
By item~1, \(\sigma S\) is not valid. For singular \(S\), the
assignment \(v\) satisfies \(S\) if and only if \(D_v\entails S\), and by
item~1 this holds if and only if \(D'_v\entails\sigma S\). Hence
\([\sigma S]=[S]\), and so \(\mathcal C(\sigma S)=\sigma\mathcal C(S)\)
with the same number of members. Also \(T\entails M_v\) if and only if
\(\sigma T\entails M'_v\), so \(K_{\sigma T}=\sigma K_T\). By
\cref{thm:p15-power-optimal-split} the two scores are the same ratio of
cardinalities.

\emph{Item 3.} The function \(r\circ\sigma\) takes values in \([0,1]\).
If \(U\) is valid, then \(\sigma U\) is valid by item~1, so
\(r(\sigma U)=1\). If \(U_1\land U_2\) is contradictory, then so is
\(\sigma U_1\land\sigma U_2=\sigma(U_1\land U_2)\). Hence
\(r(\sigma(U_1\lor U_2))=r(\sigma U_1\lor\sigma U_2)
=r(\sigma U_1)+r(\sigma U_2)\). So \(r\circ\sigma\) satisfies
\cref{def:p15-power-range-measure}. Its content measure is
\(1-r\circ\sigma=g\circ\sigma\). Since \(\sigma\) commutes with the
connectives,
\[
\begin{aligned}
 s_{r\circ\sigma}(T,S)&=\frac{g(\sigma T\lor\sigma S)}{g(\sigma S)}
 =s_r(\sigma T,\sigma S),\\[2pt]
 p_{r\circ\sigma}(T,S)&=\frac{r(\sigma T\land\sigma S)}{r(\sigma S)}
 =p_r(\sigma T,\sigma S).
\end{aligned}
\]
In each case the two denominators are the same number, so either side is
defined exactly when the other is.
\end{proof}
\end{annexitem}

\begin{annexitem}{annex:p15-w02c-p15-029-duals}{thm:p15-w02c-p15-029-duals}{Duality of minimal and maximal sentences and of the normal forms}
\begin{proof}
\emph{Item 1.} For \(v\in\Omega\), let \(\ell_i^v\) be \(\neg A_i\) if
\(v_i=1\) and \(A_i\) if \(v_i=0\). Then \(M_v=\ell_1^v\lor\cdots\lor
\ell_n^v\). The operation \(D\) leaves metavariables fixed and commutes
with negation, so \(D(\ell_i^v)=\ell_i^v\). It turns each disjunction
into a conjunction, so applying it along the bracketing of \(M_v\) gives
\(D(M_v)=\ell_1^v\land\cdots\land\ell_n^v\). This conjunction contains
\(A_i\) exactly when \(v_i=0\), that is, when \(\bar v_i=1\), and
\(\neg A_i\) exactly when \(\bar v_i=0\). By definition it is
\(D_{\bar v}\). The same computation with conjunction and disjunction
exchanged gives \(D(D_v)=M_{\bar v}\). For \(n=0\), \(M\) is the empty
disjunction \(\bot\) and the state description is the empty conjunction
\(\top\). Here \(D(\bot)=\top\) and \(D(\top)=\bot\), as the definition
of \(D\) stipulates. Since \(\bar{\bar v}=v\), the map \(v\mapsto\bar v\)
is its own inverse, so it is a bijection.

\emph{Item 2.} We prove \(D(B)\equiv\neg B(\neg A_1,\ldots,\neg A_n)\) for
every Boolean expression \(B\) by induction on \(B\). For \(B=A_i\), the
right side is \(\neg\neg A_i\equiv A_i=D(A_i)\). For \(B=\top\), the
right side is \(\neg\top\equiv\bot=D(\top)\), and dually for \(\bot\).
For \(B=\neg C\),
\(D(B)=\neg D(C)\equiv\neg\neg C(\neg A)\), which is
\(\neg B(\neg A)\). For \(B=C\land E\), De Morgan's law gives
\(D(B)=D(C)\lor D(E)\equiv\neg C(\neg A)\lor\neg E(\neg A)
\equiv\neg(C\land E)(\neg A)\). The case \(B=C\lor E\) is the same with
the connectives exchanged.

At an assignment \(w\), the sentence \(U(\neg A_1,\ldots,\neg A_n)\) has
the truth value that \(U\) has at \(\bar w\), because each atom is read
with the opposite value. Hence \(D(U)\) is true at \(w\) if and only if
\(U\) is false at \(\bar w\), that is, if and only if \(\bar w\notin[U]\).
Writing \(v=\bar w\) gives
\([D(U)]=\{\bar v:v\notin[U]\}\).

\emph{Item 3.} By item~1 and the definition of \(D\) on conjunctions,
\(D\) maps the conjunction \(\bigwedge_{v\notin[U]}M_v\) to the
disjunction \(\bigvee_{v\notin[U]}D_{\bar v}\). For \(U\) valid both
sides are empty, with \(D(\top)=\bot\). By item~2 the index set
\(\{\bar v:v\notin[U]\}\) is \([D(U)]\). So this disjunction is the
maximal normal form of \(D(U)\) given by
\cref{thm:p15-power-normal-forms}, and it is equivalent to \(D(U)\).
Since \(v\mapsto\bar v\) is injective,
\(|[D(U)]|=|\{v:v\notin[U]\}|=|\mathcal C(U)|\).

\emph{The \(L_1\) cases.} With \(A_1=P(a)\) and \(A_2=Q(a)\),
\cref{thm:p15-w02c-p15-028-l1-forms} gives the minimal normal forms
\(M_{00}\land M_{01}\) of \(P(a)\) and \(M_{10}\) of
\(\neg P(a)\lor Q(a)\). By item~1, \(D(M_{00})=D_{11}\),
\(D(M_{01})=D_{10}\) and \(D(M_{10})=D_{01}\). So the duals are
\(D_{11}\lor D_{10}\) and \(D_{01}\). The first is the maximal normal form
of \(P(a)\) in the same theorem. The second is \(\neg P(a)\land Q(a)\),
which is \(D(\neg P(a)\lor Q(a))\) by the definition of \(D\), whereas
the maximal normal form of \(\neg P(a)\lor Q(a)\) has the three members
\(D_{11},D_{01},D_{00}\). A valid sentence has the empty minimal normal
form \(\top\), and \(D(\top)=\bot\) is the empty maximal normal form of a
contradictory sentence.
\end{proof}
\end{annexitem}

\begin{annexitem}{annex:p15-w02c-p15-029-inverse}{thm:p15-w02c-p15-029-inverse}{Inverse variation and the reach of a state description}
\begin{proof}
\emph{Item 4} is proved first. Let \(\mathfrak M\) have ground
assignment \(v\), so that \(\mathfrak M\models A_i\) exactly when
\(v_i=1\). For a singular \(U\), induction on the Boolean structure of
\(U\) shows that \(\mathfrak M\models U\) if and only if \(U\) is true at
\(v\). The atomic case is the definition of \(v\), and the connective
cases are the truth tables on both sides. So \(\mathfrak M\models U\) if
and only if \(v\in[U]\). Since \([D_w]=\{w\}\),
\(\mathfrak M\models D_w\) exactly when \(w=v\).

\emph{Item 1.} Every assignment \(v\) is the ground assignment of some
structure: take the domain of the named constants, or a single object
when there are none, interpreting the atoms as \(v\) prescribes. The
language has no identity, so distinct constants may denote distinct
objects. By item~4, \(U\entails V\) implies \([U]\subseteq[V]\).
Conversely, if \([U]\subseteq[V]\) and \(\mathfrak M\models U\), then by
item~4 the ground assignment of \(\mathfrak M\) lies in \([U]\), hence in
\([V]\), and \(\mathfrak M\models V\). Taking complements in \(\Omega\),
\([U]\subseteq[V]\) holds if and only if
\(\Omega\setminus[V]\subseteq\Omega\setminus[U]\). Since
\(v\mapsto M_v\) is injective, this holds if and only if
\(\mathcal C(V)\subseteq\mathcal C(U)\). Applying this in both directions,
\(U\equiv V\) if and only if \([U]=[V]\), if and only if
\(\mathcal C(U)=\mathcal C(V)\). By the uniqueness clause of
\cref{thm:p15-power-normal-forms}, these are exactly the conditions that
the maximal normal forms coincide and that the minimal normal forms
coincide.

\emph{Item 2.} If \(U\entails V\), item~1 gives
\(|\mathcal C(U)|\ge|\mathcal C(V)|\) and \(|[U]|\le|[V]|\). Dividing by
\(2^n\) gives the two inequalities. Since
\(0\le|\mathcal C(U)|\le2^n\), we have \(0\le g_1(U)\le1\). Moreover
\(g_1(U)=0\) if and only if \(\mathcal C(U)=\varnothing\), that is,
\([U]=\Omega\). By item~1 applied with \(\top\), this holds if and only if
\(U\) is valid. Likewise \(g_1(U)=1\) if and only if \([U]=\varnothing\),
that is, \(U\) is contradictory.

\emph{Item 3.} Let \(S,S^\ast\) be the conjunctions of \(K,K^\ast\). If
\(S\equiv S^\ast\), item~1 gives \(\mathcal C(S)=\mathcal C(S^\ast)\),
so \(g_1(K)=g_1(S)=g_1(S^\ast)=g_1(K^\ast)\). If \(K\subseteq K^\ast\),
every conjunct of \(S\) is a conjunct of \(S^\ast\), so
\(S^\ast\entails S\), and item~2 gives \(g_1(S^\ast)\ge g_1(S)\).

\emph{The general measure.} Monotonicity of a range measure under
entailment is part of \cref{thm:p15-power-range-laws}. For the failure of
the converse, \cref{thm:p15-power-strength-boundary} gives a range
measure with weight \(1/2\) on each of \(00\) and \(11\), and \(0\) on
\(01\) and \(10\). Under it \(r(A)=r(B)=1/2\), but neither of
\(A\entails B\) and \(B\entails A\) holds.
\end{proof}
\end{annexitem}

\begin{annexitem}{annex:p15-w02c-p15-029-uniform-independence}{thm:p15-w02c-p15-029-uniform-independence}{Independence of disjoint atoms under the uniform measure}
\begin{proof}
Let \(I\) be the set of indices of the ground atoms occurring in \(T\),
and \(J\) the set of those occurring in \(S\), with \(I\cap J=\varnothing\).
The truth value of \(T\) at an assignment \(v\) depends only on the
restriction \(v|_I\), and that of \(S\) only on \(v|_J\). Hence there are
sets \(X\subseteq\{0,1\}^I\) and \(Y\subseteq\{0,1\}^J\) with
\([T]=\{v:v|_I\in X\}\) and \([S]=\{v:v|_J\in Y\}\). Since \(I\) and
\(J\) are disjoint, the restriction maps \(v\mapsto(v|_I,v|_J,v|_R)\),
where \(R\) is the set of remaining indices, form a bijection of
\(\Omega\) with \(\{0,1\}^I\times\{0,1\}^J\times\{0,1\}^R\). Counting,
\[
 |[T\land S]|=|X|\,|Y|\,2^{|R|},\quad
 |[T]|=|X|\,2^{|J|+|R|},\quad
 |[S]|=|Y|\,2^{|I|+|R|}.
\]
With \(n=|I|+|J|+|R|\),
\[
 r_1(T)\,r_1(S)=\frac{|X|\,|Y|\,2^{|I|+|J|+2|R|}}{2^{2n}}
 =\frac{|X|\,|Y|\,2^{|R|}}{2^{n}}=r_1(T\land S).
\]
If \(r_1(S)>0\), dividing by \(r_1(S)\) gives \(p_{r_1}(T,S)=r_1(T)\).
The sentences \(\neg T\) and \(\neg S\) contain the same atoms as \(T\)
and \(S\). So if \(g_1(S)=r_1(\neg S)>0\), then by
\cref{thm:p15-power-probability-dual}
\[
 s_{r_1}(T,S)=p_{r_1}(\neg T,\neg S)=r_1(\neg T)=g_1(T).
\]

In Carnap's example the atoms are \(P(a_1),\ldots,P(a_{101})\). The data
\(e\) and every variant \(e'\) use only the first hundred atoms, while
\(h\) uses only the last. Each of \(e,e'\) is a conjunction of literals,
true at exactly \(2\) of the \(2^{101}\) assignments. So
\(0<r_1(e')<1\), and both conditional values are defined. They equal
\(r_1(h)=1/2\) and \(g_1(h)=1/2\).

For \(r_2=m^\ast\) with names \(a,b\) and one unary \(P\): the four state
descriptions fall into three classes under exchange of \(a\) and \(b\),
namely \(\{P(a)\land P(b)\}\), \(\{P(a)\land\neg P(b),\neg P(a)\land P(b)\}\)
and \(\{\neg P(a)\land\neg P(b)\}\). By \cref{def:p15-carnap-star} each
class has mass \(1/3\), divided equally within the class. This gives the
weights \(1/3,1/6,1/6,1/3\), which are positive and sum to \(1\). So
\(r_2\) is a range measure on the singular sentences, by the
finite-weight construction of \cref{thm:p15-power-range-laws}. Then
\(r_2(P(a))=1/3+1/6=1/2\) and \(r_2(\neg P(a))=1/6+1/3=1/2\), and
\[
\begin{aligned}
 p_{r_2}(P(b),P(a))&=\frac{1/3}{1/2}=\frac23,\\[2pt]
 s_{r_2}(P(b),P(a))&=p_{r_2}(\neg P(b),\neg P(a))\\
 &=\frac{1/3}{1/2}=\frac23,
\end{aligned}
\]
using \cref{thm:p15-power-probability-dual} for the second value. Finally
\(r_2(P(b))=1/3+1/6=1/2\), so \(g_{r_2}(P(b))=1/2\).
\end{proof}
\end{annexitem}

\begin{annexitem}{annex:p15-w02c-p15-029-strict-probability}{thm:p15-w02c-p15-029-strict-probability}{Failure of an elementary law excludes a range-measure representation}
\begin{proof}
Suppose, for contradiction, that a range measure \(r\) satisfies
\(c(U,V)=p_r(U,V)\), with \(r(V)>0\), at every pair \((U,V)\) occurring
in the terms of the failing instance. Then every \(p_r\)-term of that
instance is defined. The side condition of the identity is also
satisfied, since it is the condition assumed for \(c\). This includes
\(r(T_1\land S)>0\) in the case of the product law, because
\((T_2,T_1\land S)\) is then one of the pairs. By
\cref{thm:p15-power-eight-identities} the instance holds for \(p_r\).
Replacing each \(p_r\)-term by the equal \(c\)-term, it holds for \(c\),
which contradicts the hypothesis.

For the second statement, suppose \(c(X,Y)=s_r(\neg X,\neg Y)\), with
\(g(\neg Y)>0\), at every pair \((X,Y)\) occurring in the failing
instance. By \cref{thm:p15-power-probability-dual} and invariance of
\(r\) under equivalence,
\(s_r(\neg X,\neg Y)=p_r(\neg\neg X,\neg\neg Y)=p_r(X,Y)\), and the
denominator \(r(Y)=g(\neg Y)\) is positive. Hence \(c(X,Y)=p_r(X,Y)\),
with \(r(Y)>0\), at all those pairs, which the first statement
excludes.
\end{proof}
\end{annexitem}

\begin{annexitem}{annex:p15-w02c-p15-029-extension}{thm:p15-w02c-p15-029-extension}{The generalized measure over a uniform singular part}
\begin{proof}
\emph{Item 1.} By definition \(T^\flat\) is the conjunction of the
minimal sentences \(M_v\) with \(T\entails M_v\). So \(T\entails T^\flat\),
and hence \(\neg T^\flat\entails\neg T\). It follows that
\(\neg T\land\neg T^\flat\equiv\neg T^\flat\). Therefore
\[
 \neg T\land\neg S\equiv
 (\neg T^\flat\land\neg S)\lor(\neg T\land T^\flat\land\neg S),
\]
and the two disjuncts are incompatible. By additivity and equivalence
invariance,
\[
 r(\neg T\land\neg S)=r(\neg T^\flat\land\neg S)
 +r(\neg T\land T^\flat\land\neg S).
\]
The sentences \(\neg S\) and \(\neg T^\flat\land\neg S\) are singular, so
their \(r\)-values are their \(r_1\)-values. Since \(S\) is not valid,
\(g(S)=r(\neg S)=r_1(\neg S)=g_1(S)>0\) by
\cref{thm:p15-w02c-p15-029-inverse}, so \(s_r(T,S)\) is defined. By
\cref{thm:p15-power-probability-dual} and De Morgan's law,
\[
 s_r(T,S)=\frac{r(\neg T\land\neg S)}{r(\neg S)}
 =\frac{r_1(\neg(T^\flat\lor S))}{r_1(\neg S)}
 +\frac{r(\neg T\land T^\flat\land\neg S)}{g_1(S)} .
\]
The first fraction is \(g_1(T^\flat\lor S)/g_1(S)\), which is
\(s_0(T,\{S\})\) by \cref{thm:p15-power-uniform-formula}. The second
fraction is nonnegative.

\emph{Item 2.} If \(T\) is singular, then \(T^\flat\) is the minimal
normal form of \(T\). A singular \(T\) entails \(M_v\) exactly when
\(v\notin[T]\), so \(T^\flat\equiv T\) by
\cref{thm:p15-power-normal-forms}. Then
\(\neg T\land T^\flat\land\neg S\) is contradictory, its measure is
\(0\), and item~1 gives the equality.

\emph{Item 3.} For \(v\in\{00,01,10,11\}\), let \(\mathfrak N_v\) be the
structure with domain \(\{a\}\) in which \(P(a)\) and \(Q(a)\) have the
truth values \(v_1\) and \(v_2\). Let \(\mathfrak N_{00}^+\) have domain
\(\{a,e\}\), with \(a\) in neither \(P\) nor \(Q\), and \(e\) in \(P\)
but not in \(Q\). Define
\[
 r'(U)=\tfrac14\sum_{v}\mathbf 1_{\{\mathfrak N_v\models U\}},
 \qquad
 r''(U)=\tfrac14\Bigl(\mathbf 1_{\{\mathfrak N_{00}^+\models U\}}
   +\sum_{v\neq00}\mathbf 1_{\{\mathfrak N_v\models U\}}\Bigr).
\]
Both are range measures on all sentences of \(L_1\), by the
finite-mixture construction of \cref{thm:p15-power-range-laws}. The
ground assignment of \(\mathfrak N_v\) is \(v\), and that of
\(\mathfrak N_{00}^+\) is \(00\). By item~4 of
\cref{thm:p15-w02c-p15-029-inverse}, a singular \(U\) is true in exactly
those of the four structures whose ground assignments lie in \([U]\). So
\(r'(U)=r''(U)=|[U]|/4=r_1(U)\).

The sentence \(S=P(a)\land Q(a)\) is false exactly in the structures with
ground assignment other than \(11\), so \(r'(\neg S)=r''(\neg S)=3/4\).
The theory \(T=\forall x(P(x)\to Q(x))\) is false in \(\mathfrak N_{10}\),
whose only object is in \(P\) and not in \(Q\). It is true in
\(\mathfrak N_{11}\), \(\mathfrak N_{01}\) and \(\mathfrak N_{00}\), since
no object of these is in \(P\) without being in \(Q\). It is false in
\(\mathfrak N_{00}^+\), because of the object \(e\). Hence
\(r'(\neg T\land\neg S)=1/4\), from \(\mathfrak N_{10}\), and
\(r''(\neg T\land\neg S)=2/4\), from \(\mathfrak N_{10}\) and
\(\mathfrak N_{00}^+\). By \cref{thm:p15-power-probability-dual},
\[
 s_{r'}(T,S)=\frac{1/4}{3/4}=\frac13,\qquad
 s_{r''}(T,S)=\frac{2/4}{3/4}=\frac23 .
\]
The value \(1/3\) is \(s_0(T,\{P(a),Q(a)\})\) by
\cref{thm:p15-power-source-examples}. As a check on item~1: here
\(T^\flat=\neg P(a)\lor Q(a)\), and \(\neg T\land T^\flat\land\neg S\)
holds in \(\mathfrak N_{00}^+\) and in no other of the structures. This
gives the extra term \((1/4)/(3/4)=1/3\) under \(r''\) and \(0\) under
\(r'\).

\emph{Item 4.} The sentence \(\forall x\,P(x)\) is true in
\(\mathfrak N_{11}\) and \(\mathfrak N_{10}\) and false in
\(\mathfrak N_{01}\) and \(\mathfrak N_{00}\). So
\(r'(\forall x\,P(x))=1/2\) and \(g_{r'}(\forall x\,P(x))=1/2>0\), and
\(s_{r'}(U,\forall x\,P(x))\) is defined for every sentence \(U\). The
sentence contains a variable, so it is not singular, and
\cref{def:p15-power-optimal-split} does not apply to it as data.
\end{proof}
\end{annexitem}

\begin{annexitem}{annex:p15-w02c-p15-029-converse}{thm:p15-w02c-p15-029-dualism-converse}{Converse dualism}
\begin{proof}
\emph{Item 1.} We use induction on \(B\), which is built from
metavariables, \(\top\), \(\bot\), \(\neg\), \(\land\) and \(\lor\),
implications having been expanded. For a metavariable,
\(D(D(A_i))=D(A_i)=A_i\). For constants,
\(D(D(\top))=D(\bot)=\top\), and dually for \(\bot\). For a negation,
\(D(D(\neg C))=D(\neg D(C))=\neg D(D(C))=\neg C\). For a conjunction,
\(D(D(C\land E))=D(D(C)\lor D(E))=D(D(C))\land D(D(E))=C\land E\). The
disjunction case is the same with the connectives exchanged.

\emph{The structural identity.} For every Boolean expression \(B\) and
sentences \(U_1,\ldots,U_m\),
\[
 D(B)(U)\equiv\neg B(\neg U). \tag{\(\ast\)}
\]
This is proved by the same induction as item~2 of
\cref{thm:p15-w02c-p15-029-duals}, with the sentences \(U_i\) in place of
the atoms. For \(B=A_i\) both sides are \(U_i\) and \(\neg\neg U_i\).
The constants are exchanged by \(D\) and by negation. The negation step
is immediate, and the conjunction and disjunction steps are De Morgan's
laws.

\emph{Item 2.} By \((\ast\)) and invariance of \(r\) under equivalence,
\[
 r\bigl(D(C)(U)\bigr)=r\bigl(\neg C(\neg U)\bigr)
 =1-r\bigl(C(\neg U)\bigr)=g\bigl(C(\neg U)\bigr).
\]
The left side is the denominator of
\(p_r(D(B)(U),D(C)(U))\), and the right side is the denominator of
\(s_r(B(\neg U),C(\neg U))\). Hence one term is defined exactly when the
other is. When they are defined, \cref{thm:p15-power-probability-dual}
gives
\[
 s_r\bigl(B(\neg U),C(\neg U)\bigr)
 =p_r\bigl(\neg B(\neg U),\neg C(\neg U)\bigr)
 =p_r\bigl(D(B)(U),D(C)(U)\bigr).
\]
The last step applies \((\ast)\) to \(B\) and to \(C\). It uses
equivalence invariance of \(r\) in the numerator, where
\(\neg B(\neg U)\land\neg C(\neg U)\equiv D(B)(U)\land D(C)(U)\), and in
the denominator.

\emph{Item 3.} Let \(\Psi\) be universally valid. Fix a language, a range
measure \(r\) on it and sentences \(U_1,\ldots,U_m\) at which every term
\(p_r(D(B_j)(U),D(C_j)(U))\) of \(\Psi^D\) is defined and every
hypothesis \(D(C)(U)\entails D(B)(U)\) of \(\Psi^D\) holds. Put
\(U'_i=\neg U_i\), a sentence of the same language. This is the
substitution of negated sentences for the metavariables of \(\Psi\).

\emph{Terms.} By item~2 with \(U\), each term
\(s_r(B_j(U'),C_j(U'))\) of \(\Psi\) at \(U'\) is defined and equals the
corresponding term \(p_r(D(B_j)(U),D(C_j)(U))\) of \(\Psi^D\) at \(U\).

\emph{Hypotheses.} By \((\ast)\), \(D(C)(U)\equiv\neg C(U')\) and
\(D(B)(U)\equiv\neg B(U')\). So the hypothesis of \(\Psi^D\) says
\(\neg C(U')\entails\neg B(U')\). By contraposition this is equivalent
to \(B(U')\entails C(U')\), which is the corresponding hypothesis of
\(\Psi\) at \(U'\).

\emph{Conclusion.} Since \(\Psi\) is universally valid and all its terms
are defined and all its hypotheses hold at \(U'\), \(\Psi\) is true at
\(U'\). The numbers, numerical operations, equalities and inequalities
of \(\Psi^D\) are those of \(\Psi\), and each term of \(\Psi^D\) at \(U\)
has the same value as the corresponding term of \(\Psi\) at \(U'\). So
\(\Psi^D\) is true at \(U\). The language, the measure and the sentences
were arbitrary, subject only to the definedness and hypothesis
conditions of \(\Psi^D\). Hence \(\Psi^D\) is universally valid.

\emph{Item 4.} By \cref{thm:p15-power-dualism}, the forward
transformation makes the replacements
\[
 p_r(B',C')\ \mapsto\ s_r\bigl(D(B'),D(C')\bigr),\qquad
 B'\entails C'\ \mapsto\ D(C')\entails D(B').
\]
Applied to \(\Psi^D\), it turns each
term \(p_r(D(B_j),D(C_j))\) into
\[
 s_r\bigl(D(D(B_j)),D(D(C_j))\bigr)=s_r(B_j,C_j)
\]
by item~1. It turns the hypothesis
\(D(C)\entails D(B)\) into \(D(D(B))\entails D(D(C))\), which is
\(B\entails C\). So it returns \(\Psi\). In the same way, the converse
transformation applied to the forward image of a \(p\)-statement returns
that statement. The two transformations are therefore mutually inverse.

\emph{The examples.} The power product law has the metavariables
\(T_1,T_2,S\), and \(D\) fixes each of them. Also
\(D(T_1\lor T_2)=T_1\land T_2\) and \(D(T_1\lor S)=T_1\land S\). So
\(\Psi^D\) is the displayed probability product law. By item~2 its
definedness conditions \(r(S)>0\) and \(r(T_1\land S)>0\) correspond to
\(g(S)>0\) and \(g(T_1\lor S)>0\). For the monotonicity statement, the
hypothesis \(T\entails U\) becomes \(D(U)\entails D(T)\), which is
\(U\entails T\). The conclusion \(s_r(T,S)\ge s_r(U,S)\) becomes
\(p_r(T,S)\ge p_r(U,S)\).

\emph{The warning.} For \(g(S)>0\), \(s_r(\bot,S)=g(S)/g(S)=1\), since
\(\bot\lor S\equiv S\). Replacing \(s\) by \(p\) with \(\bot\) kept gives
\(p_r(\bot,S)=r(\bot)/r(S)=0\) whenever \(r(S)>0\). The converse
transformation gives \(p_r(D(\bot),D(S))=p_r(\top,S)=r(S)/r(S)=1\). For
the copied hypothesis, \(T=\bot\) and \(U=\top\) satisfy
\(\bot\entails\top\), but \(p_r(\bot,S)=0<1=p_r(\top,S)\).
\end{proof}
\end{annexitem}

%% toc: P15 164 164.12
\section{Romeijn on inductive logic and statistics: proofs}
\label{sec:annex-p15-w02c-p15-029-romeijn}

The notation is that of \cref{sec:p15-w02c-p15-029-romeijn}. Conditional
probability is the ratio \(P(A\mid B)=P(A\cap B)/P(B)\), defined when
\(P(B)>0\). For a sequence \(s_t\), \(t_k\) is its number of entries
\(k\); \(s_0\) is the empty sequence, the whole space.

\begin{annexitem}{annex:p15-w02c-p15-029-rj-entailment}{thm:p15-w02c-p15-029-rj-entailment}{The two examples of non-ampliative entailment}
\begin{proof}
\emph{Item 1.} With \(n=3\), every \(\omega\in K^{\Omega}\) has exactly
one first entry \(\omega(1)\in\{0,1,2\}\). So \(q^0_1,q^1_1,q^2_1\) are
pairwise disjoint and their union is \(K^{\Omega}\). For a finitely
additive \(P\) with \(P(K^{\Omega})=1\),
\[
 P(q^2_1)=1-P(q^0_1)-P(q^1_1)=1-\tfrac12-\tfrac13=\tfrac16 .
\]
So every model of the premises is a model of the conclusion. For
existence, let \(\omega_k\) be the constant sequence with every entry
\(k\) and put
\(P(A)=\tfrac12[\omega_0\in A]+\tfrac13[\omega_1\in A]+\tfrac16[\omega_2\in A]\)
for every \(A\in\mathcal Q\), where \([\cdot]\) is \(1\) or \(0\) as the
condition holds or not. Each bracket is a countably additive
\(\{0,1\}\)-valued measure with total mass \(1\), and the weights sum to
\(1\), so \(P\) is a probability on \(\mathcal Q\); and
\(P(q^0_1)=\tfrac12\), \(P(q^1_1)=\tfrac13\) because \(\omega_k\in q^j_1\)
exactly when \(j=k\).

\emph{Item 2.} Since \(P(h)=a>0\), the definition of conditional
probability gives \(P(h\cap s_t)=P(h)P(s_t\mid h)=ab\). Since
\(P(s_t)=c>0\), \(P(h\mid s_t)=P(h\cap s_t)/P(s_t)=ab/c\). The argument
used only the three premises and the definition, so it holds in every
model of them.
\end{proof}
\end{annexitem}

\begin{annexitem}{annex:p15-w02c-p15-029-rj-instantial}{thm:p15-w02c-p15-029-rj-instantial}{Relevance}
\begin{proof}
\emph{Item 1.} Let \(0<\lambda<\infty\). Positivity is by induction on
\(t\): \(P(s_0)=1\), and if \(P(s_t)>0\) then (4) applies and
\[
 P(s_t\cap q^k_{t+1})=P(s_t)\,\frac{t_k+\lambda/n}{t+\lambda}>0,
\]
since \(\lambda/n>0\).

Fix \(t\ge1\), \(s_{t-1}\) with counts \(t'_j\), and \(k\). The sets
\(s_{t-1}\cap q^j_t\), \(j\in K\), partition \(s_{t-1}\). Each has
positive probability, and \(s_{t-1}\cap q^j_t\) has \(t'_k+[j=k]\)
entries \(k\) among its \(t\) entries, where \([j=k]\) is \(1\) if
\(j=k\) and \(0\) otherwise. By additivity and two applications of (4),
\begin{align*}
 \frac{P(q^k_{t+1}\cap s_{t-1})}{P(s_{t-1})}
 &=\sum_{j\in K}\frac{P(s_{t-1}\cap q^j_t)}{P(s_{t-1})}\cdot
   \frac{P(s_{t-1}\cap q^j_t\cap q^k_{t+1})}{P(s_{t-1}\cap q^j_t)}\\
 &=\sum_{j\in K}\frac{t'_j+\lambda/n}{t-1+\lambda}\cdot
   \frac{t'_k+[j=k]+\lambda/n}{t+\lambda}\\
 &=\frac{(t'_k+\lambda/n)\sum_{j}(t'_j+\lambda/n)+(t'_k+\lambda/n)}{(t-1+\lambda)(t+\lambda)} .
\end{align*}
Since \(\sum_j t'_j=t-1\) and there are \(n\) values of \(j\),
\(\sum_j(t'_j+\lambda/n)=t-1+\lambda\). The numerator is therefore
\((t'_k+\lambda/n)(t+\lambda)\), and
\(P(q^k_{t+1}\mid s_{t-1})=(t'_k+\lambda/n)/(t-1+\lambda)\). The other
side of (5) is (4) applied to \(s_{t-1}\cap q^k_t\):
\((t'_k+1+\lambda/n)/(t+\lambda)\).

Put \(a=t'_k+\lambda/n\) and \(m=t-1+\lambda\); both are positive. Then
\[
 \frac{a+1}{m+1}-\frac am=\frac{m-a}{m(m+1)},
\]
and
\[
 m-a=(t-1-t'_k)+\lambda\Bigl(1-\frac1n\Bigr)>0,
\]
because \(t'_k\le t-1\), \(\lambda>0\) and \(n\ge2\). This is the strict
inequality (5).

\emph{Item 2.} Let \(\lambda=0\). For \(t\ge1\), (4) reads
\(P(q^k_{t+1}\mid s_t)=t_k/t\). Let \(t\ge2\) and let every entry of
\(s_{t-1}\) be \(k\), with \(P(s_{t-1})>0\). Then
\(P(q^k_t\mid s_{t-1})=(t-1)/(t-1)=1\) and \(P(q^j_t\mid s_{t-1})=0\) for
\(j\ne k\). So \(P(s_{t-1}\cap q^k_t)=P(s_{t-1})>0\), and (4) applied to
the all-\(k\) sequence \(s_{t-1}\cap q^k_t\) of length \(t\) gives
\(P(q^k_{t+1}\mid s_{t-1}\cap q^k_t)=t/t=1\). For the right side,
\(P(q^k_{t+1}\cap s_{t-1}\cap q^j_t)\le P(s_{t-1}\cap q^j_t)=0\) for
\(j\ne k\), hence
\[
 P(q^k_{t+1}\cap s_{t-1})=P(q^k_{t+1}\cap s_{t-1}\cap q^k_t)=P(s_{t-1}\cap q^k_t)=P(s_{t-1}),
\]
and \(P(q^k_{t+1}\mid s_{t-1})=1\). The two sides of (5) are equal.

If instead \(P(s_{t-1})>0\) and \(s_{t-1}\) has no entry \(k\), then
\(t-1\ge1\) gives \(P(q^k_t\mid s_{t-1})=0/(t-1)=0\), so
\(P(s_{t-1}\cap q^k_t)=0\), and the conditional probability on the left
of (5) is not defined by the ratio. For \(t=0\), (4) has value
\((0+0)/(0+0)\), which is undefined.
\end{proof}
\end{annexitem}

\begin{annexitem}{annex:p15-w02c-p15-029-rj-beta}{thm:p15-w02c-p15-029-rj-beta}{Emma's farm and the \texorpdfstring{\(\lambda\)}{lambda}-rule}
\begin{proof}
\emph{A Beta identity.} For \(a,b>0\) the integral
\(B(a,b)=\int_0^1\theta^{a-1}(1-\theta)^{b-1}d\theta\) is finite and
positive: the integrand is positive on \((0,1)\), is at most
\(2^{|b-1|}\theta^{a-1}\) on \((0,\tfrac12]\) and at most
\(2^{|a-1|}(1-\theta)^{b-1}\) on \([\tfrac12,1)\), and
\(\int_0^{1/2}\theta^{a-1}d\theta=2^{-a}/a<\infty\), similarly at \(1\).
We show
\[
 B(a+1,b)=\frac{a}{a+b}B(a,b).
\]
First, \(\theta+(1-\theta)=1\) gives \(B(a+1,b)+B(a,b+1)=B(a,b)\).
Second, for \(0<\delta<\tfrac12\), integration by parts with
\(u=\theta^a\) and \(v=-(1-\theta)^b/b\) gives
\[
 \int_\delta^{1-\delta}\theta^a(1-\theta)^{b-1}d\theta
 =\Bigl[-\frac{\theta^a(1-\theta)^b}{b}\Bigr]_\delta^{1-\delta}
 +\frac ab\int_\delta^{1-\delta}\theta^{a-1}(1-\theta)^b\,d\theta .
\]
The bracket equals \(\bigl(\delta^a(1-\delta)^b-(1-\delta)^a\delta^b\bigr)/b\),
which tends to \(0\) as \(\delta\to0\) because \(a,b>0\). Letting
\(\delta\to0\), \(B(a+1,b)=\frac ab B(a,b+1)=\frac ab\bigl(B(a,b)-B(a+1,b)\bigr)\),
so \((a+b)B(a+1,b)=aB(a,b)\).

\emph{The model.} For each \(\theta\in[0,1]\), let \(P_{h_\theta}\) be
the function on the cylinder algebra with
\(P_{h_\theta}(s_t)=\theta^{t_1}(1-\theta)^{t_0}\) (with \(0^0=1\)). It
is consistent, since
\(P_{h_\theta}(s_t\cap q^1_{t+1})+P_{h_\theta}(s_t\cap q^0_{t+1})=\theta^{t_1}(1-\theta)^{t_0}(\theta+1-\theta)\),
and \(P_{h_\theta}(s_0)=1\). Let \(\pi(\theta)\) be the density (8) with
\(\mathrm{Norm}=1/B(\tfrac\lambda2,\tfrac\lambda2)\), finite because
\(\lambda>0\). The joint law assigns to \(\{\theta\in A\}\cap s_t\), for
a Borel set \(A\subseteq[0,1]\), the value
\(\int_A\pi(\theta)P_{h_\theta}(s_t)\,d\theta\); \(P_{h_\theta}\) is its
law of the observations given the parameter value \(\theta\).

\emph{Item 1.} For \(\theta\in(0,1)\) every \(P_{h_\theta}(s_{t-1})\) is
positive and
\(P_{h_\theta}(q^1_t\mid s_{t-1})=\theta^{t_1+1}(1-\theta)^{t_0}/(\theta^{t_1}(1-\theta)^{t_0})=\theta\),
where \(t_0,t_1\) are the counts of \(s_{t-1}\); likewise for \(q^0_t\)
with value \(1-\theta\). At \(\theta\in\{0,1\}\) the same holds whenever
\(P_{h_\theta}(s_{t-1})>0\). This is (7). Writing \(P_{h_\theta}(s_t)\)
as the telescoping product of the ratios
\(P_{h_\theta}(s_i)/P_{h_\theta}(s_{i-1})\) of its initial segments
gives the product in (9), each factor being \(\theta\) or \(1-\theta\)
according as \(k_i=1\) or \(0\).

\emph{Item 2.} Taking \(A=[0,1]\),
\[
 P(s_t)=\mathrm{Norm}\int_0^1\theta^{\lambda/2+t_1-1}(1-\theta)^{\lambda/2+t_0-1}d\theta
 =\frac{B(\tfrac\lambda2+t_1,\tfrac\lambda2+t_0)}{B(\tfrac\lambda2,\tfrac\lambda2)}>0 .
\]
Hence for every Borel \(A\),
\(P(\theta\in A\mid s_t)=\int_A\pi(\theta)P_{h_\theta}(s_t)\,d\theta/P(s_t)\),
which is integration over \(A\) of the density
\(\pi(\theta)P_{h_\theta}(s_t)/P(s_t)\); this is (6) in density form, and
it equals
\(\theta^{\lambda/2-1+t_1}(1-\theta)^{\lambda/2-1+t_0}/B(\tfrac\lambda2+t_1,\tfrac\lambda2+t_0)\),
which is (10).

\emph{Item 3.} By the definition of the joint law,
\[
 P(q^1_{t+1}\cap s_t)=\int_0^1\pi(\theta)\,\theta^{t_1+1}(1-\theta)^{t_0}d\theta
 =\int_0^1\theta\cdot\pi(\theta)P_{h_\theta}(s_t)\,d\theta .
\]
Dividing by \(P(s_t)\) and using item~2,
\(P(q^1_{t+1}\mid s_t)=\int_0^1\theta\,P(h_\theta\mid s_t)\,d\theta\),
which is (12); and since \(\theta=P_{h_\theta}(q^1_{t+1}\mid s_t)\) for
\(\theta\in(0,1)\) by item~1, the integrand is the one in (11), the
endpoints having Lebesgue measure zero. With
\(a=\tfrac\lambda2+t_1\), \(b=\tfrac\lambda2+t_0\),
\[
 \int_0^1\theta\,P(h_\theta\mid s_t)\,d\theta=\frac{B(a+1,b)}{B(a,b)}=\frac{a}{a+b}
 =\frac{t_1+\frac\lambda2}{t+\lambda},
\]
by the Beta identity and \(t_0+t_1=t\). This is (13).

\emph{Item 4.} For \(\lambda=2\), (13) is \((t_1+1)/(t+2)\), and
\cref{thm:p15-w02b-p15-025-ml-cstar} defines \(c^*=(s_M+w)/(s+k)\),
which with \(k=2\), \(w=1\), \(s=t\), \(s_M=t_1\) is the same number. Also
(8) is then constant, \(\theta^0(1-\theta)^0=1\), with
\(\mathrm{Norm}=1/B(1,1)=1\).

\emph{The endpoint.} For \(\lambda=0\), (8) has the shape
\(\theta^{-1}(1-\theta)^{-1}\ge\theta^{-1}\), and
\(\int_\delta^{1/2}\theta^{-1}d\theta=\log\frac1{2\delta}\to\infty\) as
\(\delta\to0\); no constant normalizes it.
\end{proof}
\end{annexitem}

\begin{annexitem}{annex:p15-w02c-p15-029-rj-np-tables}{thm:p15-w02c-p15-029-rj-np-tables}{The two tests of the farm example}
\begin{proof}
\emph{Item 1.} The likelihood ratios of \(h_1\) to \(h_0\) are
\(0.40/0\), read as \(+\infty\); \(0.30/0.05=6\); and
\(0.30/0.95=\tfrac6{19}\). For a real threshold \(r\), (14) gives
\(F(q^0)=1\) always, \(F(q^1)=1\) exactly when \(r<6\), and
\(F(q^2)=1\) exactly when \(r<\tfrac6{19}\). Test (15) has
\(F(q^0)=F(q^1)=1\) and \(F(q^2)=0\), so (14) equals (15) exactly when
\(r<6\) and \(r\ge\tfrac6{19}\). Its significance is
\(P_{h_0}(q^0)+P_{h_0}(q^1)=0+0.05\) and its power
\(P_{h_1}(q^0)+P_{h_1}(q^1)=0.40+0.30\).

The significance of a rejection region \(R\) is
\(\sum_{q^k\in R}P_{h_0}(q^k)\). Any \(R\) containing \(q^2\) has
significance at least \(0.95\). The four regions without \(q^2\) have
significance \(0\), \(0\), \(0.05\), \(0.05\) for \(\varnothing\),
\(\{q^0\}\), \(\{q^1\}\), \(\{q^0,q^1\}\), and power \(0\), \(0.40\),
\(0.30\), \(0.70\). The largest is \(0.70\), attained by (15).

\emph{Item 2.} The ratios are \(0.40/0.05=8\), \(0.30/0.05=6\) and
\(0.30/0.90=\tfrac13\). The test rejecting exactly on \(q^0\) needs
\(r<8\), \(r\ge6\) and \(r\ge\tfrac13\), that is \(6\le r<8\). Its
significance is \(P_{h'_0}(q^0)=0.05\) and its power
\(P_{h_1}(q^0)=0.40\). Under \(h'_0\), any region containing \(q^2\) has
significance at least \(0.90\), and \(\{q^0,q^1\}\) has \(0.10\). The
remaining regions \(\varnothing\), \(\{q^0\}\), \(\{q^1\}\) have
significance \(0\), \(0.05\), \(0.05\) and power \(0\), \(0.40\),
\(0.30\).

\emph{Item 3.} By items 1 and 2, the first test has \(F(q^1)=1\) and the
second has \(F(q^1)=0\). The equalities
\(P_{h_0}(q^1)=P_{h'_0}(q^1)=0.05\) are read from the tables.
\end{proof}
\end{annexitem}

\begin{annexitem}{annex:p15-w02c-p15-029-rj-interval-prior}{thm:p15-w02c-p15-029-rj-interval-prior}{Exact form of the interval-prior rendering}
\begin{proof}
The formulas for \(g_1,g_0\) are Bayes's theorem for the two-point
hypothesis space with \(P(f^1\mid h_0)=\alpha\),
\(P(f^1\mid h_1)=1-\beta\), \(P(f^0\mid h_0)=1-\alpha\),
\(P(f^0\mid h_1)=\beta\), \(P(h_0)=p\), \(P(h_1)=1-p\). Both have the
form \(g(p)=cp/(cp+d(1-p))\) with \(c,d>0\): \((c,d)=(\alpha,1-\beta)\)
for \(g_1\) and \((1-\alpha,\beta)\) for \(g_0\).

\emph{Item 1.} The denominator \(cp+d(1-p)\) is at least
\(\min(c,d)>0\) on \([0,1]\), so \(g\) is continuous, with \(g(0)=0\)
and \(g(1)=1\). For \(p<p'\),
\begin{align*}
 g(p')-g(p)&=\frac{cp'\bigl(cp+d(1-p)\bigr)-cp\bigl(cp'+d(1-p')\bigr)}{\bigl(cp'+d(1-p')\bigr)\bigl(cp+d(1-p)\bigr)}\\
 &=\frac{cd(p'-p)}{\bigl(cp'+d(1-p')\bigr)\bigl(cp+d(1-p)\bigr)}>0 .
\end{align*}
A continuous strictly increasing function maps \((l,u)\) onto
\((g(l),g(u))\) by the intermediate value theorem, and its infimum and
supremum there are \(g(l)\) and \(g(u)\), not attained.

\emph{Item 2.} \(g_1(p)<\tfrac12\) iff \(2\alpha p<\alpha p+(1-\beta)(1-p)\)
iff \(p(1-\beta+\alpha)<1-\beta\) iff \(p<b_1\), where
\(b_1=\frac{1-\beta}{1-\beta+\alpha}\). And \(g_0(p)>\tfrac12\) iff
\((1-\alpha)p>\beta(1-p)\) iff \(p>b_0\), where
\(b_0=\frac{\beta}{\beta+1-\alpha}\). Since
\(P(h_1\mid f^j)=1-P(h_0\mid f^j)\), \(P(h_0\mid f^1)<\tfrac12\) is
equivalent to \(P(h_0\mid f^1)<\tfrac12<P(h_1\mid f^1)\), and similarly for
\(f^0\). If \(l\ge b_0\) and \(u\le b_1\), every \(p\in(l,u)\) has
\(b_0<p<b_1\). Conversely, if \(l<b_0\), the nonempty interval
\((l,\min(u,b_0))\) consists of priors with \(p\le b_0\), for which
\(g_0(p)\le\tfrac12\); and if \(u>b_1\), the priors in
\((\max(l,b_1),u)\) have \(g_1(p)\ge\tfrac12\). A nonempty \((l,u)\) with
\(b_0\le l<u\le b_1\) exists iff \(b_0<b_1\), that is
\(\beta(1-\beta+\alpha)<(1-\beta)(\beta+1-\alpha)\). Expanding, the left
side is \(\beta-\beta^2+\alpha\beta\) and the right side
\(1-\alpha-\beta^2+\alpha\beta\), so the condition is
\(\beta<1-\alpha\).

\emph{Item 3.} For \(p\in(0,1)\), \(g_1(p)<p\) iff
\(\alpha<\alpha p+(1-\beta)(1-p)\) iff \(\alpha(1-p)<(1-\beta)(1-p)\) iff
\(\alpha<1-\beta\); and \(g_0(p)>p\) iff
\(1-\alpha>(1-\alpha)p+\beta(1-p)\) iff \((1-\alpha)(1-p)>\beta(1-p)\) iff
\(\beta<1-\alpha\). Both conditions are \(\alpha+\beta<1\); with
\(\alpha+\beta>1\) the same steps give the reversed inequalities. By
item~1, over \((l,u)\) with \(0<l<u<1\), \(\inf g_1=g_1(l)\) and
\(\sup g_1=g_1(u)\), which are then below \(l\) and \(u\) respectively.

\emph{Item 4.} With \(\alpha=0.05\), \(\beta=0.30\):
\(b_0=0.30/1.25=\tfrac6{25}\) and \(b_1=0.70/0.75=\tfrac{14}{15}\).
By item~2 the boundary values are \(g_0(\tfrac6{25})=\tfrac12\) and
\(g_1(\tfrac{14}{15})=\tfrac12\) (equality holds at \(b_0\), \(b_1\) by
the same equivalences with ``\(=\)''). Further,
\begin{align*}
 g_0\bigl(\tfrac{14}{15}\bigr)&=\frac{0.95\cdot14}{0.95\cdot14+0.30\cdot1}=\frac{13.3}{13.6}=\frac{133}{136},\\
 g_1\bigl(\tfrac6{25}\bigr)&=\frac{0.05\cdot6}{0.05\cdot6+0.70\cdot19}=\frac{0.3}{13.6}=\frac{3}{136}.
\end{align*}
By item~1 these are the posterior intervals. At the rounded end
\(0.93\),
\(g_0(0.93)=0.8835/0.9045\approx0.977\) and
\(g_1(0.93)=0.0465/0.0955\approx0.487\).
\end{proof}
\end{annexitem}

\begin{annexitem}{annex:p15-w02c-p15-029-rj-total-evidence}{thm:p15-w02c-p15-029-rj-total-evidence}{Coarse conditioning}
\begin{proof}
For a two-point prior \(P(\text{null})=p\), \(P(h_1)=1-p\), and an event
\(E\) with \(P(E\mid\text{null})=c>0\), \(P(E\mid h_1)=d>0\), Bayes's
theorem gives \(P(\text{null}\mid E)=g(p)=cp/(cp+d(1-p))\).

\emph{Item 1.} For \(E=q^1\), \(c=0.05\) under both \(h_0\) and
\(h'_0\), and \(d=0.30\). The posterior expressions are identical.

\emph{Item 2.} In the first problem the test (15) has
\(f^1=q^0\cup q^1\), with \(c=0+0.05\) and \(d=0.40+0.30=0.70\). In the
second, the test of \cref{thm:p15-w02c-p15-029-rj-np-tables} item~2 has
\(f^0=q^1\cup q^2\), with \(c=0.05+0.90=0.95\) and
\(d=0.30+0.30=0.60\). For fixed \(p\in(0,1)\),
\(g(p)=1/\bigl(1+(d/c)(1-p)/p\bigr)\) is strictly decreasing in the
ratio \(d/c\); the two ratios are \(14\) and \(\tfrac{12}{19}\), so the
two posteriors differ. As in the proof of
\cref{thm:p15-w02c-p15-029-rj-interval-prior} item~3, \(g(p)<p\) iff
\(c<d\) and \(g(p)>p\) iff \(c>d\); here \(0.05<0.70\) and
\(0.95>0.60\). At \(p=\tfrac12\), \(g=c/(c+d)\), giving
\(0.05/0.35=\tfrac17\), \(0.05/0.75=\tfrac1{15}\) and
\(0.95/1.55=\tfrac{19}{31}\).
\end{proof}
\end{annexitem}

\begin{annexitem}{annex:p15-w02c-p15-029-rj-mle-coverage}{thm:p15-w02c-p15-029-rj-mle-coverage}{Coverage}
\begin{proof}
By (7) and the chain rule, as in the proof of
\cref{thm:p15-w02c-p15-029-rj-beta} item~1, every \(s_t\) has
\(P_{h_\theta}(s_t)=\theta^{t_1}(1-\theta)^{t_0}\). This is the
independent Bernoulli model with success probability \(\theta\), and
the count \(t_1\) has the binomial law
\(P_{h_\theta}(t_1=j)=\binom tj\theta^j(1-\theta)^{t-j}\), because exactly
\(\binom tj\) sequences of length \(t\) have \(j\) entries \(1\).

\emph{Item 1.} For \(t=0\), \(s_0\) is the whole space and
\(P_{h_{\theta'}}(s_0)=1\) for every \(\theta'\), so the condition of
(16) holds for every \(\theta\in\Theta\).

\emph{Item 2.} The event \(t_1/t\notin[\theta-\Delta,\theta+\Delta]\)
is the event \(|t_1/t-\theta|>\Delta\). By \cref{thm:p15-fe-binomial} with
\(n=t\) and \(\epsilon=\Delta\), the probability of that event is at most
\(\theta(1-\theta)/(t\Delta^2)\). Hence
\(1\ge c_t\ge1-\theta(1-\theta)/(t\Delta^2)\to1\).

\emph{Item 3.} Let \(\theta=\tfrac12\), \(\Delta=\tfrac14\); the interval
is \([\tfrac14,\tfrac34]\) and \(P_{h_{1/2}}(t_1=j)=\binom tj2^{-t}\). For
\(t=4\), \(j/4\in[\tfrac14,\tfrac34]\) iff \(j\in\{1,2,3\}\), so
\(c_4=(4+6+4)/16=\tfrac78\). For \(t=5\), \(j/5\in[\tfrac14,\tfrac34]\)
iff \(\tfrac54\le j\le\tfrac{15}4\) iff \(j\in\{2,3\}\), so
\(c_5=(10+10)/32=\tfrac58<\tfrac78\).
\end{proof}
\end{annexitem}

\begin{annexitem}{annex:p15-w02c-p15-029-rj-lambda-zero}{thm:p15-w02c-p15-029-rj-lambda-zero}{The \texorpdfstring{\(\lambda=0\)}{lambda = 0} rule and Fisher's estimate}
\begin{proof}
\emph{Item 1.} With \(\lambda=0\) and \(n=2\), (4) is
\((t_1+0)/(t+0)=t_1/t\) for \(t\ge1\). Under (7) the likelihood of
\(h_\theta\) on \(s_t\) is (9), \(\theta^{t_1}(1-\theta)^{t_0}\), and by
\cref{thm:p15-fe-binomial} its unique maximizer on the closed interval
\([0,1]\) is \(t_1/t\); so (16) is \(\{t_1/t\}\). For \(s^{000101}\),
\(t=6\) and \(t_1=2\), giving \(\tfrac13\).

\emph{Item 2.} \(P(q^0_1)>0\), and (4) gives
\(P(q^0_2\mid q^0_1)=1/1\) and \(P(q^0_3\mid s^{00})=2/2\), each
conditioning event having positive probability by the preceding step.
So \(P(s^{000})=P(q^0_1)\cdot1\cdot1>0\). Then
\(P(q^1_4\mid s^{000})=0/3=0\) and
\(P(s^{0001})=P(s^{000})\cdot0=0\). The ratio
\(P(q^0_5\cap s^{0001})/P(s^{0001})\) has denominator \(0\) and is
undefined. For \(s^{0001}\), \(t=4\), \(t_1=1\), so
\(\hat\theta=\tfrac14\) and
\(P_{h_{1/4}}(q^0_5)=1-\tfrac14=\tfrac34\) by (7).
\end{proof}
\end{annexitem}

\begin{annexitem}{annex:p15-w02c-p15-029-rj-fiducial-transfer}{thm:p15-w02c-p15-029-rj-fiducial-transfer}{Transfer (24)}
\begin{proof}
Write \(h_\theta\) for the event \(\{\theta\}\times W\), \(\omega\) for
\(\Theta\times\{\omega\}\), and \(\{\hat\theta=x\}\) for
\(\{(\theta,\omega):f(\theta,\omega)=x\}\). By the uniqueness assumption,
\(f(\theta,\omega)=x\) holds iff \(\theta=f^{-1}(x,\omega)\).

\emph{Item 1.} If \(P(h_\theta)>0\), then by (23)
\[
 P(\hat\theta=x\mid h_\theta)=\frac{\sum_{\omega:\,f(\theta,\omega)=x}P(h_\theta)P(\omega)}{P(h_\theta)}
 =P(\{\omega:f(\theta,\omega)=x\}).
\]

\emph{Item 2.} The event \(h_\theta\cap\{\hat\theta=x\}\) is
\(\{(\theta,\omega):f(\theta,\omega)=x\}=\{(\theta,\omega):f^{-1}(x,\omega)=\theta\}\),
and the event \(\{\hat\theta=x\}\) is
\(\{(f^{-1}(x,\omega),\omega):\omega\in W\}\). So, within
\(\{\hat\theta=x\}\), \(h_\theta\) is the union of the events \(\omega\)
with \(f^{-1}(x,\omega)=\theta\). Hence
\begin{align*}
 P(h_\theta\mid\hat\theta=x)&=\sum_{\omega:\,f^{-1}(x,\omega)=\theta}P(\omega\mid\hat\theta=x)\\
 &=\sum_{\omega:\,f^{-1}(x,\omega)=\theta}P(\omega)=P(\{\omega:f^{-1}(x,\omega)=\theta\}),
\end{align*}
using the premise \(P(\omega\mid\hat\theta=x)=P(\omega)\).

\emph{Item 3.} Here \(f(\theta,\omega)=x\) iff \(\theta=x-\omega\bmod2\),
so \(f\) is invertible as required. Put
\(P(h_\theta,\omega)=P(h_\theta)P(\omega)\) with \(P(h_0)=1\), \(P(h_1)=0\)
and \(P(\omega)=\tfrac12\); (23) holds by construction. Almost surely
\(\theta=0\), so \(\hat\theta=\omega\) and \(P(\hat\theta=x)=\tfrac12>0\).
Then \(P(h_0\mid\hat\theta=x)=P(h_0\cap\{\hat\theta=x\})/\tfrac12=\tfrac12/\tfrac12=1\),
while \(P(\{\omega:x-\omega\equiv0\})=P(\omega=x)=\tfrac12\). The premise of
item~2 fails as well: \(P(\omega=x\mid\hat\theta=x)=1\ne\tfrac12\).

\emph{Item 4.} Let the prior be \(\pi\) and the error law \(\rho\) on
\(\mathbb Z_m\). By (23),
\(P(\hat\theta=x,\omega=w)=\pi(x-w)\rho(w)\), and
\(P(\hat\theta=x)=\sum_w\pi(x-w)\rho(w)\).

If \(\pi\equiv1/m\), then \(P(\hat\theta=x)=1/m>0\) and
\(P(\omega=w\mid\hat\theta=x)=\rho(w)\), the premise of item~2. Item~2
then gives \(P(h_\theta\mid\hat\theta=x)=P(\{\omega:x-\omega=\theta\})=\rho(x-\theta)\).
By Bayes's theorem and item~1,
\(P(h_\theta\mid\hat\theta=x)=\pi(\theta)P(\hat\theta=x\mid h_\theta)/P(\hat\theta=x)=(1/m)\rho(x-\theta)/(1/m)\),
the same value.

If \(\rho\equiv1/m\), then \(P(\hat\theta=x)=\frac1m\sum_w\pi(x-w)=\frac1m\),
because \(x-w\) runs over all of \(\mathbb Z_m\) as \(w\) does. So
\(P(\omega=w\mid\hat\theta=x)=\pi(x-w)\). This equals \(\rho(w)=1/m\) for
every \(w\) iff \(\pi(\theta)=1/m\) for every \(\theta\), again because
\(x-w\) runs over all of \(\mathbb Z_m\). Thus the premise holds for one
\(x\), and then for every \(x\), exactly when \(\pi\) is uniform.
\end{proof}
\end{annexitem}
